\section{Appendix: Key Commands and Environment Variables}
\label{sec:AppendixOperations}

This appendix gathers the most frequently used commands and configuration variables mentioned elsewhere in the manual.

\subsection{Core commands}

\begin{center}
\begin{tabular}{|p{.28\textwidth}|p{.62\textwidth}|}
\hline
\textbf{Area} & \textbf{Typical command} \\
\hline
Java unit tests & \shellcmd{mvn -f implementations/java/pom.xml -q test} \\
\hline
Java integration tests & \shellcmd{mvn -f implementations/java/pom.xml -pl it verify} \\
\hline
PostgreSQL test DB setup & \shellcmd{python3 shared/db/postgresql/setup-for-test.py} \\
\hline
DB2 test DB setup & \shellcmd{python3 shared/db/db2/setup-for-test.py} \\
\hline
Erlang compile & \texttt{rebar3 compile} \\
\hline
Erlang tests & \texttt{rebar3 eunit} \\
\hline
Rust tests & \texttt{cargo test} \\
\hline
Rust Python extension & \texttt{maturin develop} \\
\hline
Rust Python tests & \shellcmd{python3 -m pytest -q python_tests/test_ipto.py} \\
\hline
\end{tabular}
\end{center}

\subsection{Important PostgreSQL variables}

The PostgreSQL-backed implementations reuse a common naming scheme for local development:
\begin{itemize}
\item \texttt{IPTO\_PG\_HOST}
\item \texttt{IPTO\_PG\_PORT}
\item \texttt{IPTO\_PG\_USER}
\item \texttt{IPTO\_PG\_PASSWORD}
\item \texttt{IPTO\_PG\_DATABASE}
\end{itemize}

In Java, the default JDBC configuration is instead read from \texttt{configuration\_pg.xml}, but the practical target is the same local PostgreSQL instance used by tests and development tooling.

\subsection{Important Neo4j variables}

The Erlang and Rust implementations both use Neo4j-related environment variables. The common ones are:
\begin{itemize}
\item \texttt{IPTO\_NEO4J\_URL}
\item \texttt{IPTO\_NEO4J\_DATABASE}
\item \texttt{IPTO\_NEO4J\_USER}
\item \texttt{IPTO\_NEO4J\_PASSWORD}
\end{itemize}

Additional tuning variables exist in the Erlang implementation for timeouts, retries and bootstrap behavior.

\subsection{GraphQL SDL-related settings}

Where the GraphQL layer is configurable at runtime, the most important settings are the ones that determine where SDL is read from. In the Erlang implementation, for example, this includes:
\begin{itemize}
\item app env \texttt{graphql\_schema\_sdl},
\item app env \texttt{graphql\_schema\_file},
\item env var \texttt{IPTO\_GRAPHQL\_SCHEMA\_FILE}.
\end{itemize}

In the Java implementation, the concrete application packaging currently provides a GraphQL schema resource under \filepath{implementations/java/quarkus-app/src/main/resources/ipto.graphqls}.

\subsection{Practical advice}

When working on IPTO, the most reliable sequence is usually:
\begin{enumerate}
\item provision the database or backend environment first,
\item run the tests closest to the subsystem being changed,
\item if SDL or persistence semantics changed, run an integration-level check rather than only unit tests,
\item only then verify other implementation surfaces affected by the same conceptual change.
\end{enumerate}

This is a direct consequence of the architecture described in the manual: many changes cross the boundaries between schema, catalog, runtime and language binding.

\subsection{Example unit payload}

The repository-facing JSON format used by the current implementations is generic. A minimal unit payload looks like this:

\begin{lstlisting}[language=Java,
    framesep=8pt,
    xleftmargin=5pt,
    framexleftmargin=5pt,
    frame=tb,
    framerule=0pt]
{
  "tenantid": 1,
  "corrid": "550e8400-e29b-41d4-a716-446655440000",
  "status": 30,
  "unitname": "example-unit",
  "attributes": [
    {
      "attrid": 101,
      "attrname": "dcterms:title",
      "attrtype": 1,
      "value": ["Example title"]
    }
  ]
}
\end{lstlisting}

The important points are:
\begin{itemize}
\item the unit header identifies tenant, correlation id, status and optional name,
\item each attribute occurrence carries the repository attribute id and datatype id,
\item primitive values are represented as arrays, even when a single value is present,
\item nested records follow the same idea recursively, but use record-valued attributes and child attribute arrays.
\end{itemize}

This is the payload shape that is translated into \texttt{repo\_unit\_kernel}, \texttt{repo\_unit\_version}, \texttt{repo\_attribute\_value} and the vector tables during ingestion.
